\appendixsection{主要系统调用表}\label{sec:appendix-syscalls}

本附录列出当前系统中与用户态 Shell 和常用程序运行关系较密切的系统调用。不同调用在内核中的具体处理函数分布于进程、文件、时间、信号和终端等模块，表中仅保留论文正文分析所需的接口层信息。

\begin{table}[htbp]
  \caption{当前系统主要系统调用}\label{tab:appendix-syscalls}
  \centering
  \begin{tabular}{p{0.2\textwidth}p{0.22\textwidth}p{0.46\textwidth}}
    \toprule
    \textbf{类别} & \textbf{接口} & \textbf{作用} \\
    \midrule
    进程控制 & \verb|fork|、\verb|exec|、\verb|wait|、\verb|exit| & 创建子进程、替换进程图像、等待子进程退出和结束当前进程 \\
    文件访问 & \verb|open|、\verb|creat|、\verb|read|、\verb|write|、\verb|close| & 访问普通文件、设备文件和管道文件对象 \\
    目录与链接 & \verb|chdir|、\verb|stat|、\verb|link|、\verb|unlink|、\verb|mkdir| & 维护目录项、路径解析结果和文件元数据 \\
    描述符管理 & \verb|dup|、\verb|pipe| & 支撑 Shell 管道、重定向和父子进程间通信 \\
    信号与时间 & \verb|signal|、\verb|kill|、\verb|sigret|、\verb|sleep| & 安装处理函数、发送信号、从信号处理返回和定时睡眠 \\
    \bottomrule
  \end{tabular}
\end{table}

\appendixsection{地址空间与磁盘布局}\label{sec:appendix-layout}

当前 RISC-V 内核由 QEMU 加载到 \verb|0x80200000|，物理内存基址为 \verb|0x80000000|。内核建立高地址映射后，通过 \verb|0xc0000000| 起始的窗口访问物理内存，并通过 \verb|0xf0000000| 起始的 MMIO 映射访问 UART、PLIC 和 VirtIO 等设备。用户空间从低地址开始布置，ELF 加载器将可装入段、数据区和用户栈映射到进程页表中；\verb|exec| 完成后，\verb|argc| 和 \verb|argv| 位于用户栈顶附近。

\begin{table}[htbp]
  \caption{当前系统主要布局约定}\label{tab:appendix-layout}
  \centering
  \begin{tabular}{p{0.28\textwidth}p{0.26\textwidth}p{0.36\textwidth}}
    \toprule
    \textbf{对象} & \textbf{位置或格式} & \textbf{说明} \\
    \midrule
    DRAM 物理基址 & \verb|0x80000000| & QEMU virt 平台内存起始地址 \\
    内核加载地址 & \verb|0x80200000| & OpenSBI 转交后内核入口所在区域 \\
    内核高地址窗口 & \verb|0xc0000000| 起 & 内核访问物理内存和进程用户区的主要窗口 \\
    MMIO 映射窗口 & \verb|0xf0000000| 起 & UART、PLIC、VirtIO 等设备访问区域 \\
    用户程序格式 & ELF64/RISC-V & 由用户态 C 运行时和交叉编译工具链生成 \\
    磁盘镜像 & \verb|target/disk.img| & 由构建工具写入 Shell、用户程序和文件系统内容 \\
    \bottomrule
  \end{tabular}
\end{table}

\appendixsection{QEMU 启动和调试命令}\label{sec:appendix-qemu}

当前系统的常用运行入口由 Makefile 提供。普通交互运行使用 \verb|make qemu-riscv64|，该目标启动 \verb|qemu-system-riscv64| 的 virt 机器，使用 OpenSBI 默认固件，并通过 VirtIO MMIO 设备挂载文件系统镜像。需要调试内核时可使用 \verb|make qemu-riscv64g| 启动 GDB stub，再由交叉 GDB 连接到对应端口。Rust 内核的类型检查可通过 \verb|make check| 执行，构建过程会使用 RISC-V 目标和裸机标准库子集。

\appendixsection{关键数据结构说明}\label{sec:appendix-structures}

\begin{table}[htbp]
  \caption{正文涉及的关键数据结构}\label{tab:appendix-structures}
  \centering
  \begin{tabular}{p{0.26\textwidth}p{0.22\textwidth}p{0.42\textwidth}}
    \toprule
    \textbf{结构或类型} & \textbf{所在模块} & \textbf{作用} \\
    \midrule
    \verb|TrapContext| & \verb|src/interrupt/| & 保存 Trap 入口需要恢复的寄存器与控制状态 \\
    \verb|TaskContext| & \verb|src/proc/| & 保存内核态上下文切换所需的寄存器状态 \\
    \verb|Process| & \verb|src/proc/| & 表达进程控制块、状态、信号和用户区关联关系 \\
    \verb|Userspace| & \verb|src/user/| & 保存进程私有的文件描述符、当前目录、参数和内存描述 \\
    \verb|KernelPages|、\verb|UserPages| & \verb|src/mm/| & 表达内核页和用户页资源的所有权边界 \\
    \verb|InodeRef|、\verb|FileRef| & \verb|src/fs/| & 封装 inode 与打开文件对象的引用计数和释放路径 \\
    \bottomrule
  \end{tabular}
\end{table}
